\section{Kimi K2.5：三个维度的融合}

\subsection{训练与架构}

K2.5 将上述三个维度的进展融合为一个模型：
\begin{itemize}
    \item \textbf{Muon + QK-Clip 优化器}：提升 token efficiency
    \item \textbf{Kimi Delta Attention + Linear 架构}：增强 long context 能力
    \item \textbf{Agent Swarms}：开辟并行 scaling 新维度
\end{itemize}

模型在 NVIDIA H800 GPU 上训练，每个节点包含 2 TB RAM，GPU 间通过 NVLink 互联。K2.5 基础模型经过超过 30 万亿 token 的训练（base model 15T + K2.5 额外 15T），训练过程极其平稳——没有任何 loss spike。

\begin{importantbox}{Early Fusion：原生视觉-文本联合训练}
K2.5 是首个具有\textbf{原生联合视觉-文本能力}的开放模型。不同于在文本模型之上添加视觉能力的 late fusion 方案，K2.5 从训练第 0 天起就融合视觉和文本 token（early fusion），初步实验表明这一方案优于 late fusion。
\end{importantbox}

\subsection{跨模态增强效应}

K2.5 训练中发现了一个令人兴奋的现象：两种模态可以相互增强。

\begin{importantbox}{Vision 改善 Text，Text 改善 Vision}
\begin{itemize}
    \item \textbf{Vision $\rightarrow$ Text}：仅使用视觉任务进行 RL（不含任何数学/编程任务），却能提升文本推理性能
    \item \textbf{Text $\rightarrow$ Vision}：拥有强大的文本基础后，完全不需要视觉 SFT 数据（Zero Vision SFT），仅用文本 SFT + 联合 RL 就能达到接近 SOTA 的视觉性能
\end{itemize}
这一发现表明，early fusion 使两种模态真正共享了表示空间，实现了能力的双向迁移。
\end{importantbox}


